% !TEX root = ../../../main.tex
% Sources: LEC02, IMG_9122--IMG_9125; LEC03, IMG_9126--IMG_9128.
\section{The completeness proof chain}
\label{sec:analysis-i:completeness-equivalences}

The preceding section derives the Archimedean property directly from
LUBP and proves $1/n\to0$ and $2^{-n}\to0$. Thus $\R$ has the
Archimedean hypothesis needed for the shrinking intervals in this chain.
In any Archimedean ordered field the same estimates hold.

\begin{theorem}
  In an Archimedean ordered field, the six completeness properties in
  \autoref{sec:analysis-i:supremum} are equivalent.
\end{theorem}
We follow the lecture's cycle in \autoref{fig:completeness-cycle}.
The final implication closes the argument at the beginning of week three.
The Archimedean hypothesis ensures that the bisection lengths used below
actually tend to zero.
\begin{figure*}[htbp]
  \centering
  \begin{tikzpicture}[analysis diagram]
  \node[analysis property] (lub) at (0,0) {LUBP};
  \node[analysis property] (mct) at (2.4,0) {MCT};
  \node[analysis property] (nip) at (4.8,0) {NIP};
  \node[analysis property] (bwt) at (7.2,0) {BWT};
  \node[analysis property] (cc) at (9.6,0) {CC};
  \node[analysis property] (dc) at (12,0) {DC};
  \draw[->] (lub) -- (mct);
  \draw[->] (mct) -- (nip);
  \draw[->] (nip) -- (bwt);
  \draw[->] (bwt) -- (cc);
  \draw[->] (cc) -- (dc);
  \draw[->] (dc.south) -- (12,-0.9) -- (0,-0.9) -- (lub.south);
\end{tikzpicture}

  \caption[The six completeness implications.]{The six implications. Order properties appear at the ends;
    the middle steps pass through sequences and nested intervals.}
  \label{fig:completeness-cycle}
\end{figure*}

\subsection{LUBP implies MCT}
\begin{proof}
  Let $(x_n)$ be nondecreasing and bounded above, and set
  $L=\sup\{x_n\mid n\in\N\}$. For $\varepsilon>0$, the number
  $L-\varepsilon$ is not an upper bound. There is therefore an $N$ with
  $x_N>L-\varepsilon$. For $n\geq N$, monotonicity gives
  \[
    L-\varepsilon<x_N\leq x_n\leq L<L+\varepsilon.
  \]
  Thus $x_n\to L$. Applying this argument to $(-x_n)$ proves convergence
  of every nonincreasing sequence bounded below.
\end{proof}

\subsection{MCT implies NIP}
\begin{proof}
  Write the nested intervals as $I_n=[a_n,b_n]$. Their endpoints satisfy
  \[
    a_1\leq a_n\leq a_{n+1}\leq b_{n+1}\leq b_n\leq b_1.
  \]
  By MCT, $a_n\to a$ and $b_n\to b$. Passing to limits gives $a\leq b$.
  For fixed $n$, all later endpoints lie in $[a_n,b_n]$, so
  $a_n\leq a\leq b\leq b_n$. Hence $[a,b]$ is contained in every $I_n$.
  Conversely, if $x\in I_n$ for all $n$, passage to the limit in
  $a_n\leq x\leq b_n$ gives $a\leq x\leq b$. Therefore
  \[
    \bigcap_{n=1}^{\infty}I_n=[a,b]\ne\varnothing.
  \]
\end{proof}
The intersection need not be a single point. If additionally
$b_n-a_n\to0$, then $a=b$ and the intersection is a singleton.

\subsection{NIP implies BWT}
\begin{proof}
  Let $(x_n)$ be bounded, with every term in a closed interval $I_1$ of
  positive length. Bisect $I_1$. At least one closed half contains $x_n$
  for infinitely many indices $n$. Choose such a half as $I_2$, and repeat.
  This constructs nested closed intervals $I_k$ of lengths
  $|I_1|/2^{k-1}$, each containing infinitely many indexed terms.
  By NIP, choose $L\in\bigcap_k I_k$.

  Choose $n_1$ with $x_{n_1}\in I_1$. Having chosen $n_{k-1}$, choose
  $n_k>n_{k-1}$ with $x_{n_k}\in I_k$, which is possible because infinitely
  many indices qualify. Then
  \[
    |x_{n_k}-L|\leq |I_1|\,2^{1-k}\longrightarrow0.
  \]
  Thus $(x_{n_k})$ is a convergent subsequence.
\end{proof}
\autoref{fig:nested-bisection} shows a possible sequence of choices.
\begin{marginfigure}
  \centering
  \begin{tikzpicture}[interval diagram,x=8mm]
  \foreach \a/\b/\row/\name in {0/4/0/I_1,0/2/-1/I_2,1/2/-2/I_3,1/1.5/-3/I_4} {
    \node[left] at (-0.3,\row) {$\name$};
    \draw (\a,\row) -- (\b,\row);
    \node[interval closed] at (\a,\row) {};
    \node[interval closed] at (\b,\row) {};
  }
  \node at (1.25,-4) {$\vdots$};
\end{tikzpicture}

  \caption[Nested intervals under bisection.]{Schematic bisection. Each chosen interval contains infinitely
    many terms, counted by index. The process continues beyond the rows shown.}
  \label{fig:nested-bisection}
\end{marginfigure}

\subsection{BWT implies CC}
First note that convergence always implies the Cauchy condition. If
$x_n\to L$, choose $N$ so that $|x_n-L|<\varepsilon/2$ for $n\geq N$.
For $m,n\geq N$, the triangle inequality gives
\[
  |x_n-x_m|\leq|x_n-L|+|x_m-L|<\varepsilon.
\]
The converse needs completeness.

\begin{proof}
  Let $(x_n)$ be Cauchy. Choose $N_0$ so that $|x_n-x_{N_0}|<1$ for
  every $n\geq N_0$. The tail is bounded by $|x_{N_0}|+1$, and the finitely
  many earlier terms are bounded as well. For example, the maximum of
  $|x_1|,\ldots,|x_{N_0-1}|,|x_{N_0}|+1$ bounds the whole sequence.
  BWT therefore gives a subsequence $x_{n_k}\to L$.

  Given $\varepsilon>0$, choose $N$ so that $|x_n-x_m|<\varepsilon/2$
  for $m,n\geq N$. Choose a subsequence index $n_k\geq N$ with
  $|x_{n_k}-L|<\varepsilon/2$. For every $n\geq N$,
  \[
    |x_n-L|\leq|x_n-x_{n_k}|+|x_{n_k}-L|<\varepsilon.
  \]
  Hence the full sequence converges to $L$.
\end{proof}

\subsection{CC implies DC}
\begin{proof}
  Let $A,B$ be a partition as in DC, and choose $a_1\in A$, $b_1\in B$.
  At step $n$, take the midpoint $m_n=(a_n+b_n)/2$. If $m_n\in A$,
  set $a_{n+1}=m_n$ and $b_{n+1}=b_n$; otherwise set $a_{n+1}=a_n$ and
  $b_{n+1}=m_n$. The intervals are nested, their endpoints remain on
  the prescribed sides, and
  \[
    b_n-a_n=\frac{b_1-a_1}{2^{n-1}}\longrightarrow0.
  \]
  For $m\geq n$, both $a_m$ and $a_n$ lie in $[a_n,b_n]$, so
  $|a_m-a_n|\leq b_n-a_n$. The same holds for the right endpoints.
  Consequently both sequences are Cauchy. By CC they converge, and the
  vanishing difference shows that they share a limit $c$.

  Every $a\in A$ is less than every $b_n$, so $a\leq c$ on taking limits.
  Likewise every $b\in B$ is greater than every $a_n$, so $c\leq b$.
  This is the required boundary.
\end{proof}

\subsection{DC implies LUBP}
\begin{proof}
  Let $S$ be nonempty and bounded above. Let $B$ be the set of all upper
  bounds of $S$, and let $A$ be its complement in the field. By assumption,
  $B$ is nonempty. If $s\in S$, then $s-1\in A$, so $A$ is nonempty too.
  If $a\in A$, there is an $s\in S$ with $a<s$. Thus, for every $b\in B$,
  $a<s\leq b$. The partition satisfies DC, which supplies a number $c$
  with $a\leq c\leq b$ for all $a\in A$, $b\in B$.

  We claim that $c\in B$. Otherwise some $s_0\in S$ satisfies $c<s_0$.
  Its midpoint with $c$ satisfies
  \[
    c<\frac{c+s_0}{2}<s_0.
  \]
  This midpoint cannot belong to $A$, since it exceeds $c$, and cannot
  belong to $B$, since it is smaller than the element $s_0$ of $S$.
  That contradicts the partition. Hence $c$ is an upper bound of $S$.
  Since $c\leq b$ for every upper bound $b$, it is the least upper bound.\marginnote{The
  midpoint is the useful choice. Adding a fixed amount such as $1$ to $c$
  need not leave the result below $s_0$.}
\end{proof}

\subsection{Exercise on the greatest lower bound property}
The corresponding lower-bound argument can be made directly from DC,
without passing through the least upper bound property.

% !TEX root = ../../hw03.tex
% Source: HW3 Intro to Math Analysis I (1).pdf, page 1, question 1.
\begin{exercise}
  \label{ex:hw03:dedekind-infimum}
  \solutionlink{hw03:dedekind-infimum}
  Use the existence of the Dedekind cut in $\R$ to prove the greatest lower
  bound property. That is, if $S$ is a nonempty subset of $\R$ which is
  bounded below, then there is a unique number $c\in\R$ such that
  \[
    c=\inf S.
  \]
  (Do not use LUBP.)
\end{exercise}


\subsection{Exercises on the nested interval hypotheses}
Closedness and boundedness each do work in the nested interval property.
The following examples isolate these hypotheses. The Archimedean property
from \autoref{thm:analysis-i:archimedean} supplies the indices needed
to exclude any fixed point from the intersection.

% !TEX root = ../../hw02.tex
% Homework 02, Exercise 1. Shared problem statement.
\begin{exercise}
  \label{ex:hw02:nested-bounded-intervals}
  \solutionlink{hw02:nested-bounded-intervals}
  Give an example of a sequence of nested nonempty bounded intervals
  $\{I_n\}$ in $\R$ such that $I_{n+1}\subset I_n$ for each
  $n=1,2,\ldots$ but
  \[
    \bigcap_{n=1}^{\infty} I_n=\varnothing.
  \]
  Explain what goes ``wrong''.
\end{exercise}

% !TEX root = ../../hw02.tex
% Homework 02, Exercise 2. Shared problem statement.
\begin{exercise}
  \label{ex:hw02:nested-closed-intervals}
  \solutionlink{hw02:nested-closed-intervals}
  Give an example of a sequence of nested nonempty closed intervals
  $\{I_n\}$ in $\R$ such that $I_{n+1}\subset I_n$ for each
  $n=1,2,\ldots$ but
  \[
    \bigcap_{n=1}^{\infty} I_n=\varnothing.
  \]
  Explain what goes ``wrong''.
\end{exercise}

